%% =====================================================================
%% COMPSCI 760 -- Group 3 Final Report (IEEE two-column template)
%% ---------------------------------------------------------------------
%% * Compile with pdfLaTeX + BibTeX from the report/ folder.
%% * \ph{...} = red placeholder. Search for "\ph{" before submitting.
%%   Set \showplaceholdersfalse to hide them and check the page count.
%% * Limit: 12 pages EXCLUDING Contributions and References. No appendix.
%% * Figures are read from ../figures/ (the project's figures folder, written
%%   by the remake_*.ipynb notebooks). A missing file draws a red box instead.
%% * References live in references.bib (same folder). IEEEtran.bst
%%   numbers them in order of first citation. Minimum 9, max 1 page.
%% * Numbers come from: remake_dataset.ipynb (data), remake_ablation.ipynb
%%   (RQ1), remake_train.ipynb (RF/XGBoost/CatBoost, SHAP, explanations),
%%   remake_nn.ipynb, remake_traindition.ipynb (hedonic),
%%   remake_feature_reduction.ipynb (correlated features).
%% =====================================================================

\documentclass[conference]{IEEEtran}

%% ---------------------------- Packages -------------------------------
\usepackage{cite}
\usepackage{amsmath,amssymb,amsfonts}
\usepackage{graphicx}
\usepackage{xcolor}
\usepackage{booktabs}
\usepackage{multirow}
\usepackage{url}
\usepackage[hidelinks]{hyperref}

%% ---------------------- Placeholder helpers --------------------------
\newif\ifshowplaceholders
\showplaceholderstrue   % change to \showplaceholdersfalse to hide all \ph{}
\ifshowplaceholders
  \newcommand{\ph}[1]{\textcolor{red}{[#1]}}
\else
  \newcommand{\ph}[1]{}
\fi

% \figph{file}{height}{description}: includes the file if it exists,
% otherwise draws a red placeholder box of the given height.
\newcommand{\figph}[3]{%
  \IfFileExists{#1}%
    {\includegraphics[width=\linewidth]{#1}}%
    {\fbox{\parbox[c][#2][c]{0.94\linewidth}{\centering\small
      \textcolor{red}{FIGURE PLACEHOLDER\\[2pt]\texttt{\detokenize{#1}}\\[2pt]#3}}}}}

\graphicspath{{../figures/}}

\begin{document}

%% ============================ TITLE ==================================
% Rubric (22 pts, shared with Abstract): state the method and the finding.
\title{Modeling Auckland Property Prices with Explainable AI}

\author{
\IEEEauthorblockN{Dominik Wecke, Huy Nguyen, Jiayi Du, Thao Nguyen, Xiang Yang}
\IEEEauthorblockA{ University of Auckland,
Auckland, New Zealand\\
    }
}

\maketitle

%% =========================== ABSTRACT ================================
% One paragraph, 150--200 words: motivation, data, method, results, insight.
\begin{abstract}
House prices affect lending, rates and major financial decisions, yet
free online valuation tools give buyers a number without a reason. We
build an explainable machine-learning model for Auckland house prices
and test whether Auckland-specific external data (location,
neighbourhood census and deprivation statistics, flood and aircraft-noise
exposure, and development activity) improves prediction beyond a
property's own characteristics. We link 59,708 Auckland sales (January
2024 to August 2026) to map coordinates and evaluate every model with
five-fold, time-ordered cross-validation, so later sales never leak into
earlier training folds. Adding 16 external features to a fixed XGBoost
model lowers test RMSE from \$343,718 to \$212,678 and MAPE from 19.3\%
to 11.0\%, a statistically significant gain, and a two-stage residual
model shows that the external features explain 29.3\% of the error a
property-only model leaves. Permutation importance, SHAP and drop-group
retraining agree that location carries most of the gain, while
neighbourhood statistics largely repeat what location already says, and
flood, noise and development features add almost nothing. XGBoost is the
most accurate of five model families (MAPE 11.0\%, against 15.7\% for a
hedonic regression), and SHAP turns each prediction into a short,
per-sale explanation. A comparison with free online valuation estimates
is still in progress.
\end{abstract}

\begin{IEEEkeywords}
house price prediction, explainable AI, SHAP, gradient boosting,
hedonic regression, Auckland
\end{IEEEkeywords}

%% ========================= I. INTRODUCTION ===========================
% About 1--1.5 pages. Rubric: Problem statement (22 pts).

\section{Introduction}\label{sec:intro}

For most New Zealand households, a home is the largest thing they will
ever buy, and its value affects many decisions. Buyers and sellers use
valuations to set prices, banks use them to decide how much to lend, and
Auckland Council uses them to set rates. In Auckland, these valuations
are hard to get right. Homes vary widely in size, age and condition, and
the region itself varies too: some areas are close to the city centre or
the coast, some neighbourhoods are far wealthier than others, and some
are exposed to flooding, aircraft noise or fast-changing development. A
model that only looks at the house itself is therefore likely to miss a
large part of its price.

Free online valuation tools give buyers a number but not a reason. A
buyer cannot see whether an estimate is high because of the house, its
location, or a flood risk the tool has picked up. This matters when the
estimate is used to make a large financial decision.

Past research has looked at this problem from two directions. Hedonic
studies of Auckland have measured the effect of single external factors
on price, such as flood risk~\cite{samarasinghe2010}, school
zones~\cite{rehm2008}, leaky-building stigma~\cite{rehm2009} and social
housing~\cite{sequeira2021}. These models are easy to interpret, but
each studies one factor at a time, and they are built to explain effects
rather than to predict new sales~\cite{shmueli2010}. Machine learning
models predict prices more accurately than hedonic
regression~\cite{ceh2018,an2025}, and SHAP~\cite{lundberg2017} can show
which features drive their predictions. However, we found no study that
applies explainable machine learning to Auckland, and few that test how
much a combined set of external features improves prediction for sales
the model has not seen.

This paper addresses the following research questions:
\begin{itemize}
  \item \textbf{RQ1:} Does adding Auckland-specific external features
  (location, neighbourhood, hazards and development) improve house price
  prediction compared with using only a property's own features?
  \item \textbf{RQ2:} How accurate and interpretable is our explainable
  model compared with a traditional hedonic model and free online
  valuation estimates for the same Auckland properties?
\end{itemize}

To answer these questions, we built a dataset of 59,708 single-home
Auckland sales from January 2024 to August 2026. We linked each sale to
a map coordinate using probabilistic record linkage, and used the
coordinate to add 16 external features. We then trained the same XGBoost
model twice, once with and once without the external features, using
time-ordered cross-validation so that the model always learns from
earlier sales and is tested on later ones. A paired bootstrap tests
whether the difference is real, and a two-stage model tests whether
external features explain errors that property features cannot. Three
attribution methods show which groups of external features matter. We
then compare XGBoost with random forest, CatBoost, a neural network and
a hedonic regression on the same folds, and use SHAP to explain
individual predictions.

Our main contributions are:
\begin{enumerate}
  \item \textbf{A linked Auckland dataset.} We join council sales records
  to LINZ addresses using probabilistic record linkage, and add location,
  2023 Census, deprivation, hazard and development features to each sale.
  We describe each cleaning step so that the dataset can be rebuilt.
  \item \textbf{A fair test of external data.} By keeping the model and
  its settings fixed and training only on past sales, we isolate the
  effect of the external features. They reduce RMSE by 38.1\% (from
  \$343,718 to \$212,678) and MAPE from 19.3\% to 11.0\%, a significant
  improvement, and on their own they explain 29.3\% of the error left by
  a property-only model.
  \item \textbf{Evidence on what drives the gain.} Location is the main
  source of the gain. Neighbourhood statistics carry real signal but
  overlap heavily with location, and hazard and development features add
  almost nothing at city scale.
  \item \textbf{A comparison of model families with per-sale
  explanations.} XGBoost is the most accurate of five model families
  trained on identical inputs, with a MAPE of 11.0\% against 15.7\% for a
  hedonic regression. All five models agree on
  the most important features, and SHAP turns each prediction into a
  short explanation a buyer can read. \ph{Online valuation benchmark
  (RQ2): pending.}
\end{enumerate}

The rest of this paper is organised as follows.
Section~\ref{sec:related} reviews related work.
Section~\ref{sec:data} describes the data and how it was cleaned and
linked. Section~\ref{sec:method} explains the experimental design.
Section~\ref{sec:results} presents and discusses the results.
Section~\ref{sec:limits} sets out the limitations, and
Section~\ref{sec:conclusion} concludes and suggests future work.

%% ========================= II. RELATED WORK ==========================
% HARD LIMIT: 1 page. Do not copy from the literature survey.
% End each subsection with how our work differs. Rubric: 11 pts.
\section{Related Work}\label{sec:related}

\subsection{Hedonic models and external factors in Auckland}
Hedonic pricing treats a house as a bundle of characteristics, each with
its own implicit price, which can be estimated by regressing sale price
on those
characteristics~\cite{rosen1974}. Hedonic models are mainly built to
\emph{explain} the effect of each factor, not to predict new sales as
accurately as possible~\cite{shmueli2010}. Auckland studies using this
approach show that factors outside the property matter. Adding
valuer-defined submarkets to a hedonic model improved prediction more
than replacing it with more advanced spatial
models~\cite{bourassa2007}. Flood-prone homes on the North Shore sold for
about 6.2\% less before flood maps were published and 2.3\% less
afterwards~\cite{samarasinghe2010}. Being in the zone of a popular school
raised prices~\cite{rehm2008}, while leaky-building stigma lowered house
prices by about 5\%~\cite{rehm2009}, and nearby social housing lowered
prices by 1.7--23\%~\cite{sequeira2021}. Studies outside New Zealand also
show that ignoring the link between nearby sales can overstate the
effect of an external factor such as flooding~\cite{bui2024}.

These studies have two limits. First, each one looks at a single
external factor, so they do not show how much a combined set of external
data improves prediction. Second, the models are linear, and the form of
each effect (for example, how price changes with distance) must be
chosen by hand. Our work instead adds 16 external features to one model
at the same time and measures their combined effect on prediction error
for sales the model has not seen.

\subsection{Machine learning versus hedonic models}
Tree-based machine learning models learn non-linear effects and
interactions directly from the data. Across many studies they predict
prices more accurately than hedonic regression. A random forest reached
a test MAPE of about 7\% against 17\% for regression on Ljubljana
apartments~\cite{ceh2018}, and an $R^2$ of 0.94 against 0.65 on Busan
condominiums~\cite{an2025}. In Russia, a random forest also beat
regression (MAPE 14.9\% vs 18.3\%), although its error was much higher
for larger apartments~\cite{antipov2012}, so a single overall error can
hide weak spots. Gradient-boosted trees such as
XGBoost~\cite{chen2016xgboost} perform similarly well. Soltani
\textit{et al.}~\cite{soltani2022} improved tree models further by adding
a feature built from nearby recent sales, and stressed that such
features must use only sales that happened \emph{before} the sale being
predicted.

Most of these comparisons change the model, the features, or both at
once, so it is hard to tell how much of the gain comes from better data
and how much from a better algorithm. We keep the algorithm and its
settings fixed and change only the feature set, so any difference comes
from the external data. When we compare algorithms, we do the opposite:
every model gets the same features, folds and preprocessing. We also
train only on past sales and test on later ones, which avoids leaking
future information into the model.

\subsection{External data and explainable models}
Machine learning models are harder to interpret than hedonic models,
which matters when buyers and lenders need to know why a home is valued
as it is. SHAP~\cite{lundberg2017} addresses this by splitting each
prediction into the contribution of each feature. In Lisbon, adding open
data on the surroundings reduced MAE by 8.24\%, and SHAP showed that
longitude and latitude were among the most important
features~\cite{neves2024}. In Sydney, SHAP showed that access to amenities
affected price in a non-linear way that changed with distance and type
of amenity~\cite{lee2025}. In Beijing, location was the most important
factor, and the effects of other features varied between
areas~\cite{chen2025beijing}. SHAP, however, describes how the model
uses each feature, not what causes prices to change, and it cannot show
on its own that a model predicts well.

These studies use SHAP to rank features within one model. When
features carry overlapping information, as coordinates and
location-based features do, a single ranking can be misleading. We
therefore first test whether external features improve prediction, and
only then explain the model. To see how much each group of external
features matters, we use three methods (permutation importance, SHAP and
drop-group retraining) that react differently to overlapping features,
and we confirm with a separate residual test that the external features
add information the property features do not.

\subsection{Summary}
To our knowledge, no published study has applied explainable machine
learning to Auckland house prices, and few have tested the joint
predictive value of external features in a single model. Our work
addresses both gaps. It also compares the model with four other model
families, including a hedonic model, on identical folds.

%% ===================== III. DATA ====================================
\section{Data}\label{sec:data}

\subsection{Data sources}
Our main data is the Auckland Council sales file~\cite{akl_dvr}. It
lists 92,010 property sales from January 2024 to August 2026, with 41
columns that describe the sale (price, date, sale type) and the property
(land area, floor area, building age, condition, garages, and so on).

We added external data from nine sources: Statistical Area 2 (SA2)
boundaries~\cite{statsnz_sa2} and 2023 Census totals for each
SA2~\cite{statsnz_census}; the NZDep2023 deprivation
index~\cite{nzdep2023}; new dwellings consented in each SA2 from Stats
NZ~\cite{statsnz}; flood plain maps and the aircraft noise overlay from
Auckland Council~\cite{akl_gis}; the New Zealand Schools
Directory~\cite{moe_schools}; the LINZ coastline~\cite{linz_coast};
supermarket and road data from OpenStreetMap~\cite{osm}; and bus stop
locations from Auckland Transport~\cite{at_gtfs}.

\subsection{Linking sales to coordinates}\label{sec:linkage}
% ------------------------------------------------------------------
% The \ph{} items below are the Splink settings that are not recorded in
% this repository (only the output, dvrs_LINZ_matched.csv, is). Dominik:
% please fill them in from the linkage code. Every number outside a \ph{}
% was checked against dvrs_LINZ_matched.csv and the sales file.
% ------------------------------------------------------------------
The sales file describes each property by its council valuation
reference and a street address: a house number, an optional unit or
letter, and a street name with an abbreviated road type (for example
``St Michaels AV''). It has no coordinates and no suburb. To place each
property on the map, we linked it to the LINZ NZ Addresses
register~\cite{linz_addresses}, which gives a coordinate for every
official address in New Zealand. The two sources share no key, and the
same address is often written differently (abbreviations, unit formats,
missing letters or typing errors), so an exact join on the address text
would miss many properties and could not tell a near-miss from a wrong
address.

We therefore used probabilistic record linkage with
Splink~\cite{linacre2022,splink_docs}, an open-source implementation of
the Fellegi--Sunter model~\cite{fellegi1969}. For a candidate pair of
records, the model compares several fields. Each comparison has a small
number of ordered levels, such as exact agreement, close agreement under
a string-similarity measure, disagreement, and missing. For level $\ell$
of comparison $k$, the model learns two probabilities:
$m_{k\ell}$, the probability of seeing that level if the two records are
the same address, and $u_{k\ell}$, the probability of seeing it if they
are not. The evidence from all comparisons is combined into a match
weight
\begin{equation}
  w = \log_2\frac{\lambda}{1-\lambda}
      + \sum_{k}\log_2\frac{m_{k\gamma_k}}{u_{k\gamma_k}},
  \qquad
  p = \frac{2^{w}}{1+2^{w}},
  \label{eq:linkage}
\end{equation}
where $\gamma_k$ is the level observed for comparison $k$, $\lambda$ is
the prior probability that a random pair is a match, and $p$ is the
resulting match probability. Agreement on a field adds weight, and
disagreement subtracts it. The more unusual the agreement is among
non-matching pairs, the more weight it adds.

No hand-labelled pairs are needed to fit the model. The $u$
probabilities are estimated from randomly sampled pairs, almost all of
which are non-matches, and the $m$ probabilities are estimated with the
expectation--maximisation (EM) algorithm~\cite{linacre2022}. Splink can
also adjust for term frequency, so that agreeing on a rare street name
counts as stronger evidence than agreeing on a common one. To keep the
number of pairs manageable, only pairs that pass a blocking rule are
scored. \ph{Blocking rules used, e.g.\ same street name, or same house
number and first letters of the street name.} \ph{Fields compared and
their levels, e.g.\ house number (exact), unit/letter (exact or
missing), street name (exact, Jaro--Winkler $\ge$ 0.9, else), road type;
any cleaning applied first, such as lower-casing and expanding road-type
abbreviations; whether term-frequency adjustment was used.}

Each valuation reference was assigned the LINZ address with the highest
match weight, provided that the weight exceeded a threshold of
\ph{threshold, as a match weight or match probability}. \ph{Clerical
review: how many pairs were checked by hand, in which match-weight
bands, and the false-match rate found.} The result links 81,363 of the
85,180 distinct properties in the sales file (95.5\%) to a LINZ address
and coordinate. Each property receives one address. A small number of
LINZ addresses (72) are shared by two valuation references, as when two
cross-leased homes share one street number. The coordinates are joined
back to the sales by valuation reference, which combines the roll
number, the assessment number and the lower-cased suffix that separates
units sharing one assessment. After the filters in
Table~\ref{tab:funnel}, 59,798 of the 61,244 arm's-length single-home
sales (97.6\%) have a coordinate.

Once each property had a coordinate, we placed it in its 2025 SA2 to
attach the census and deprivation values, and on the council maps to
find whether it lies in a flood plain and which aircraft noise category
applies. Distances to the city centre (Auckland CBD), the nearest open
school and the coastline are straight-line distances in metres (NZTM).
Walking distances to the nearest supermarket and bus stop were measured
along the road network.

\subsection{Cleaning}
We kept only sales that represent a normal market price for one home.
Table~\ref{tab:funnel} shows the filters, applied in order.

\begin{table}[t]
  \centering
  \caption{Filters applied to the sales data, in order}
  \label{tab:funnel}
  \begin{tabular}{@{}lr@{}}
    \toprule
    Filter & Sales left \\
    \midrule
    All parsed sales                       & 92,010 \\
    Residential, rural or lifestyle zoning & 80,070 \\
    One unit of use                        & 77,370 \\
    Freehold                               & 76,395 \\
    One property per sale                  & 74,112 \\
    Arm's-length market sale               & 61,244 \\
    Coordinate found                       & 59,798 \\
    Exact duplicate records removed        & 59,711 \\
    Floor-area typos removed               & 59,708 \\
    \bottomrule
  \end{tabular}
\end{table}

The zoning and unit filters remove commercial properties and buildings
with several homes. The tenure filter removes leasehold sales, which sell
for less because the buyer does not own the land. The one-property filter
matters most: when several properties are sold together, the sale records
one combined price but the features of only one property, so the model
cannot learn from these rows. The arm's-length filter removes sales
between related parties, which may not reflect market value. We then
removed sales with no coordinate, records that repeat the same property,
date and price, and three records whose building floor area is more than
ten times the council's mass-appraisal floor area (for example
360,000~m\textsuperscript{2} for an ordinary house), which are data-entry
errors.

We also removed five columns that would leak the answer to the model. The
net, chattels and other price columns add up to the sale price. The
council's capital value and land value are themselves price estimates, so
a model that uses them would copy the council's valuation instead of
learning from property and location data.

\subsection{Features}
We built two feature sets (Table~\ref{tab:features}). The
\emph{baseline} set has 11 features that describe only the property and
the sale date; it contains no information about location. The
\emph{augmented} set adds 16 external features, for 27 in total. All SA2
statistics are area-level values shared by every home in the same SA2.

\begin{table}[t]
  \centering
  \caption{Feature sets}
  \label{tab:features}
  \begin{tabular}{@{}p{0.25\linewidth}p{0.67\linewidth}@{}}
    \toprule
    Group & Features \\
    \midrule
    \multicolumn{2}{@{}l}{\textit{Baseline (11)}} \\
    Size & land area, building floor area \\
    Building & age, condition, construction type, property type,
      off-street parking, garages under main roof, freestanding
      garages \\
    Time & sale year, sale month \\
    \midrule
    \multicolumn{2}{@{}l}{\textit{External (16)}} \\
    Location (7) & latitude, longitude, distance to CBD, distance to
      coastline, distance to nearest school, walking distance to
      supermarket, walking distance to bus stop \\
    Neighbourhood (6) & NZDep2023 index, median personal income,
      employment rate, unemployment rate, population, population
      density (all per SA2) \\
    Hazards (2) & in flood plain, aircraft noise category \\
    Development (1) & new dwellings consented in the SA2 (past 12
      months) \\
    \bottomrule
  \end{tabular}
\end{table}

Land area, floor area and the two walking distances are right-skewed,
so we applied $\log(1+x)$ to them. Categories are converted to integer
codes, and missing numbers (walking distances for 1.4--2.7\% of sales) are
filled with the median. Both steps are learned from the training data
only (Section~\ref{sec:split}).

\textbf{Removing near-duplicate features.} Two features were dropped
before modelling. Sale quarter is calculated from sale month, and the
council's mass-appraisal floor area repeats the building floor area
(Spearman $\rho = 0.96$); both pairs had variance inflation factors
(VIF) of 17--18. Dropping them cost XGBoost \$212 of RMSE (0.1\%), and
afterwards no feature has a VIF above 7.3. Removing them matters for
interpretation more than accuracy: two near-copies of a feature split
its SHAP credit between them and make both look less important than
they are.

%% ===================== IV. METHODOLOGY ==============================
\section{Methodology}\label{sec:method}
All code is available online.\footnote{\ph{https://github.com/ORG/REPO}.
The dataset is built in \texttt{remake\_dataset.ipynb}; the RQ1
experiments are in \texttt{remake\_ablation.ipynb}; the model comparison,
SHAP and per-sale explanations are in \texttt{remake\_train.ipynb},
\texttt{remake\_nn.ipynb} and \texttt{remake\_traindition.ipynb}.}

\subsection{Models}
Our main model is XGBoost~\cite{chen2016xgboost}, a gradient-boosted
tree model that predicts the sale price in dollars. We chose it because
it handles non-linear effects and mixed feature types well, and tree
models have outperformed linear regression in earlier house price
studies~\cite{ceh2018,an2025}. It uses 500 trees, maximum depth 6,
learning rate 0.05, 80\% of rows and 80\% of features sampled per tree,
and random seed 42.

To check that our conclusions do not depend on one algorithm, we train
four more models on exactly the same inputs:
\begin{itemize}
  \item \textbf{Random forest:} 150 trees, maximum depth 18, at least 2
  sales per leaf.
  \item \textbf{CatBoost:} 500 trees, depth 6, learning rate 0.05, set to
  match XGBoost. It receives the same integer-coded categories as the
  other models.
  \item \textbf{Neural network:} a multilayer perceptron with a learned
  embedding for each categorical feature, three hidden layers (256, 128
  and 64 units, ReLU, dropout 0.1), trained with Adam (learning rate
  $10^{-3}$, batch 512) on standardised log price. Training stops when
  the error on the latest 10\% of the training period has not improved
  for 10 epochs. Inputs are clipped to the 0.5th--99.5th percentiles of
  the training data, and predictions are converted back to dollars with
  Duan's smearing factor (Section~\ref{sec:hedonic}).
  \item \textbf{Hedonic regression:} see Section~\ref{sec:hedonic}.
\end{itemize}
No model is tuned. If each model were tuned separately, part of any
difference between them could come from tuning rather than from the
features or the algorithm. Fixed settings keep the comparisons fair, at
the cost of slightly understating what each model could reach.

\subsection{Training and testing}\label{sec:split}
House prices change over time, so a model should only learn from sales
that happened before the sales it predicts. Randomly shuffling the data
would let the model learn from future sales, which makes results look
better than they would be in practice~\cite{soltani2022}.

We therefore sort the 59,708 sales by date and use five-fold
expanding-window cross-validation (Fig.~\ref{fig:folds}). The data is cut
into six consecutive time blocks of about 9,950 sales. In fold $k$, the
model trains on blocks $1$ to $k$ and is tested on block $k+1$. Each fold
therefore trains on more history than the one before. Every sale after
the first block is tested exactly once, by a model that has never seen
it. We do not keep a separate validation set because no settings are
tuned.

Within each fold, all preprocessing is learned from the training blocks
only: the category codes, the median values used to fill gaps, and the
price limits described next. The test block is then processed with these
training values.

Before training, we remove sales priced below the 1st percentile or above
the 99th percentile of the training prices (\$430,000--\$450,000 and
\$4.2M--\$4.5M, depending on the fold). The same limits are applied to the test block. This stops a
few extreme sales from dominating the error. It also means our results
describe the middle 98\% of the market, not the very cheapest or most
expensive homes. After this step, 48,777 test predictions remain across
the five folds.

\begin{figure}[t]
  \centering
  \figph{../figures/timeseries_split.png}{3cm}{Expanding-window folds.}
  \caption{Five-fold expanding-window cross-validation. Each fold trains
  on all earlier sales and tests on the next period.}
  \label{fig:folds}
\end{figure}

\subsection{Evaluation metrics}
For $n$ test sales with actual price $y_i$ and predicted price
$\hat{y}_i$, we report three error measures:
\begin{align}
  \mathrm{MAE}  &= \frac{1}{n}\sum_{i=1}^{n}\lvert y_i-\hat{y}_i\rvert,
  \label{eq:mae}\\
  \mathrm{RMSE} &= \sqrt{\frac{1}{n}\sum_{i=1}^{n}(y_i-\hat{y}_i)^2},
  \label{eq:rmse}\\
  \mathrm{MAPE} &= \frac{100}{n}\sum_{i=1}^{n}
                   \frac{\lvert y_i-\hat{y}_i\rvert}{y_i}.
  \label{eq:mape}
\end{align}
MAE is the typical error in dollars. RMSE gives more weight to large
errors, which matter most for expensive homes. MAPE is the error as a
percentage of price, so a \$100K error on a \$600K home counts more than
the same error on a \$3M home. We combine the test predictions from all
five folds and compute each measure once on the combined set.

\subsection{Testing whether the improvement is real}
To check that the augmented model is better and not just lucky, we use a
paired bootstrap. We draw $n$ test sales at random with replacement and
compute the difference
\begin{equation}
  \Delta^{(b)} = M_{\text{base}}^{(b)} - M_{\text{aug}}^{(b)}
  \label{eq:bootstrap}
\end{equation}
for each metric $M$. Both models are scored on the same resampled sales,
so the comparison is fair. We repeat this $B = 2{,}000$ times. The 2.5th
and 97.5th percentiles of $\Delta^{(b)}$ give a 95\% confidence interval.
If the interval does not include zero, the improvement is statistically
significant.

\subsection{Two-stage residual test}
The test above shows whether external features help, but the augmented
model could still use them to learn something the property features
already contain. The two-stage test checks whether external features add
\emph{new} information.

In stage~1, we take the baseline model's out-of-fold test predictions
and compute each sale's error $r_i = y_i - \hat{y}_i^{\text{base}}$. In
stage~2, a new XGBoost model uses only the 16 external features to
predict $r_i$. Stage~2 has no access to property features, so it can
only succeed if location and the other external data explain what the
property features missed.

Stage~2 uses its own five-fold expanding-window split over the 48,777
sales that have a stage-1 error, again sorted by date. The first block
(8,132 sales) is used only for training, leaving 40,645 sales with a
stage-2 prediction. We report the out-of-fold $R^2$ of stage~2 in
Section~\ref{sec:results}. A value above zero means the external features
carry information the property features do not.

\subsection{Which external features matter}\label{sec:attribution}
For this analysis we use fold~5: the model trains on all sales up to
January 2026 and is tested on the 9,769 sales from January to August
2026. We study the four external groups in Table~\ref{tab:features}, and
use three methods, because each one can be misleading on its own when
features are correlated:
\begin{itemize}
  \item \textbf{Grouped permutation importance.} We randomly shuffle all
  features of a group together in the test set and measure how much the
  RMSE increases. We repeat this 20 times and report the average.
  \item \textbf{SHAP.} We compute SHAP values~\cite{lundberg2017} on
  500 test sales. Each SHAP value is the amount, in dollars, that a
  feature moves one prediction up or down. We report each group's share
  of the total mean absolute SHAP value.
  \item \textbf{Drop-group retraining.} For each group (and each single
  external feature), we retrain the model without it and measure the
  change in test RMSE.
\end{itemize}
These methods react differently when two features carry similar
information. Permutation importance breaks a group while the model still
relies on it, so it shows how much the model \emph{uses} the group, which
can overstate how much it is \emph{needed}. Drop-group retraining lets
the model fall back on a similar feature in another group, which shows
what is unique to the group but can understate its importance. SHAP
splits the credit between similar features. Reading the three together
shows both how much a group is used and how much of it is unique.

\subsection{Hedonic baseline}\label{sec:hedonic}
To compare with the traditional approach, we fit a hedonic model: a
linear regression on the logarithm of price,
\begin{equation}
  \ln y_i = \beta_0 + \mathbf{x}_i^{\top}\boldsymbol{\beta}
            + \varepsilon_i,
  \label{eq:hedonic}
\end{equation}
estimated by ordinary least squares. It uses the same 27 features as the
other models, so any difference in accuracy comes from the model, not the
inputs. Categories are one-hot encoded with the most common level as the
reference, and any column that is constant in a fold's training data is
dropped for that fold. As for the neural network, numeric inputs are
clipped to the training data's 0.5th--99.5th percentiles, so that a
single unusual value cannot produce an extreme price through the
exponential.

The model predicts log price, and turning this back into dollars with
$\exp(\cdot)$ underestimates the average price. We correct this with
Duan's smearing factor~\cite{duan1983}:
\begin{equation}
  \hat{y}_i = \exp\!\left(\hat{\beta}_0 + \mathbf{x}_i^{\top}
  \hat{\boldsymbol{\beta}}\right)\cdot
  \frac{1}{m}\sum_{k=1}^{m}e^{\hat{\varepsilon}_k},
  \label{eq:smearing}
\end{equation}
where $\hat{\varepsilon}_k$ are the $m$ training residuals. The hedonic
model uses the same folds, price limits and metrics as the other models,
so all results are directly comparable.

\subsection{Explaining individual predictions}\label{sec:explain}
SHAP values add up: for each sale, the model's average prediction plus
the SHAP value of every feature equals the prediction. We use this to
explain single predictions. For a held-out sale, we start from the
average predicted price and add each feature's contribution in dollars.
For a short summary, we sum the contributions within five themes (the
four external groups plus ``house and land''), which is valid because
SHAP values are additive. A waterfall chart shows the largest individual
contributions in real units (square metres, kilometres, deciles), so a
reader can see which characteristics raised or lowered the price. For
tree models we use exact TreeSHAP; for the neural network and hedonic
model we use SHAP's permutation explainer with 100 training sales as
background.

\subsection{Online valuation benchmark}\label{sec:onlineval}
% ------------------------------------------------------------------
% PLACEHOLDER: data collection in progress (see cache/_scratch/
% online_valuation_progress.md). Once complete, replace the \ph{} below
% with the matched sample size, the match-confirmation rule (exact sale
% date and price against the public listing/estimate history) and any
% properties excluded because no public estimate could be found.
% ------------------------------------------------------------------
To test the second half of RQ2, we draw a sample of held-out,
test-period sales and look up a free public valuation estimate for each
one (e.g.\ homes.co.nz, propertyvalue.co.nz or trademe.co.nz), matched to
the council record by address and confirmed against the actual sale date
and price so we are comparing the estimate that was available
\emph{before} the sale, not one updated afterwards using the sale itself.
We then compute the same RMSE/MAE/MAPE as the other models on this
matched subset. \ph{Sample size, match-confirmation results and the
online-valuation error metrics.}

%% ===================== V. EXPERIMENTS AND RESULTS ====================
\section{Experiments and Results}\label{sec:results}

\subsection{Does external data help?}\label{sec:res-ablation}
Table~\ref{tab:ablation} gives the pooled out-of-fold error over all five
folds. Adding the 16 external features cuts RMSE by \$131,040, from
\$343,718 to \$212,678 (a 38.1\% relative reduction), MAE by \$95,242,
and MAPE by 8.4 percentage points, from 19.3\% to 11.0\%
(Fig.~\ref{fig:ablation}). The paired bootstrap
(Eq.~\ref{eq:bootstrap}, $B=2{,}000$) gives 95\% confidence intervals of
\$126,988--\$135,157 for the RMSE improvement, \$93,358--\$97,194 for MAE
and 8.2--8.5 points for MAPE. None includes zero, so the gain is real
and not a result of how the test sales happen to fall. The augmented
model is also better in every one of the five folds, with RMSE between
\$191,483 and \$238,373 against \$320,862 to \$370,069 for the baseline.
In practical terms, an MAE of \$128,000 is still a large amount of money,
but it is close to half the \$223,000 typical error of a model that only
knows the house.

\begin{table}[t]
  \centering
  \caption{Baseline (property only) vs.\ augmented (with external
  features), XGBoost, pooled out-of-fold over 5 folds, $n=48{,}777$}
  \label{tab:ablation}
  \begin{tabular}{@{}lrrrr@{}}
    \toprule
    Model & Features & RMSE & MAE & MAPE \\
    \midrule
    Baseline  & 11 & \$343,718 & \$223,224 & 19.3\% \\
    Augmented & 27 & \$212,678 & \$127,983 & 11.0\% \\
    \bottomrule
  \end{tabular}
\end{table}

\begin{figure}[t]
  \centering
  \figph{../figures/baseline_vs_augmented_metrics.png}{3.2cm}{Baseline
  vs.\ augmented RMSE/MAE/MAPE bar chart.}
  \caption{Pooled out-of-fold error of the baseline and augmented
  XGBoost models.}
  \label{fig:ablation}
\end{figure}

\subsection{Do external features add new information?}
A model could look like it needs external data only because the
property features are weak, not because location genuinely carries new
information. The two-stage test rules this out: a model using
\emph{only} the 16 external features explains 29.3\% of the variance in
what the property-only model gets wrong ($R^2=0.293$ on 40,645 held-out
residuals), cutting the residual RMSE from \$339,582 to \$285,308. Since
this second-stage model never sees a single property feature, this
29.3\% must be information the property features do not already
contain.

\subsection{Which external features matter?}\label{sec:res-attribution}
Table~\ref{tab:attribution} compares the three attribution methods from
Section~\ref{sec:attribution} on fold~5 (9,769 test sales, augmented
RMSE \$191,483). Together, the external features are 52.3\% of the
model's total mean absolute SHAP value, and dropping all of them raises
RMSE by \$129,379, in line with the five-fold result above.

\begin{table}[t]
  \centering
  \caption{Attribution of the external feature groups on fold~5.
  Permutation: RMSE increase when the group is shuffled. SHAP: share of
  total mean $|$SHAP$|$. Drop-group: RMSE change when the model is
  retrained without the group (negative = marginally better without
  it).}
  \label{tab:attribution}
  \begin{tabular}{@{}lrrr@{}}
    \toprule
    Group & Permutation & SHAP & Drop-group \\
    \midrule
    All external (16)  & +171.2\% & 52.3\% & +\$129,379 \\
    Location (7)       & +131.1\% & 35.1\% & +\$20,146 \\
    Neighbourhood (6)  & +32.1\%  & 16.1\% & +\$1,665 \\
    Hazards (2)        & +0.3\%   & 0.4\%  & +\$370 \\
    Development (1)    & +0.6\%   & 0.6\%  & $-$\$1,489 \\
    \bottomrule
  \end{tabular}
\end{table}

Location is the most important group on all three measures. Distance to
the CBD alone is 15.5\% of total SHAP, more than latitude (8.1\%) and
longitude (8.6\%) individually, and second only to building floor area
(21.0\%) across all 27 features (Fig.~\ref{fig:shap}).

The neighbourhood group shows why the three methods must be read
together. The model clearly uses it: shuffling it raises RMSE by 32.1\%,
and it holds 16.1\% of SHAP, led by the deprivation index (6.9\%).
Yet retraining without it costs only \$1,665, because the model can
recover nearly the same information from location. The overlap runs both
ways but is not symmetric: without the location group, the model still
loses \$20,146 even though the neighbourhood statistics remain. Without
all external features it loses \$129,379, more than six times the sum
of the four single-group losses. In other words, location and neighbourhood largely
describe the same thing (how desirable an area is), and location
describes it more completely. The same pattern holds for single
features: no external feature costs more than \$3,752 when dropped on its
own (latitude), and for eight of them the model is marginally better
without them, within noise.

Hazard and development features matter little at city scale. Only 6.4\%
of sales lie in a flood plain and 5.3\% in an aircraft noise overlay, so
these
features can matter for those homes but cannot explain much of the
variation in price across Auckland, and the 12-month count of consents in
an SA2 changes slowly and overlaps with location.

\begin{figure}[t]
  \centering
  \figph{../figures/shap_xgboost_fold5.png}{4.2cm}{SHAP summary plot
  for the augmented XGBoost model.}
  \caption{SHAP summary for the augmented XGBoost model on 500 fold-5
  test sales (15 most important features). Each dot is one sale; colour
  is the feature's value.}
  \label{fig:shap}
\end{figure}

\subsection{Comparison of model families}\label{sec:res-models}
Table~\ref{tab:models} compares the five models on the same 48,777 test
sales and the same 27 features. XGBoost is the most accurate on every
metric (MAPE 11.0\%), followed closely by CatBoost and random forest
(11.6\%), then the neural network (12.3\%). The hedonic regression is
clearly the least accurate (15.7\%). All four machine-learning models
train all five folds in under a minute on a CPU.

\begin{table}[t]
  \centering
  \caption{All models on the same features and folds, pooled
  out-of-fold, $n=48{,}777$. Time is for training all five folds.}
  \label{tab:models}
  \begin{tabular}{@{}lrrrr@{}}
    \toprule
    Model & RMSE & MAE & MAPE & Time \\
    \midrule
    XGBoost             & \$212,678 & \$127,983 & 10.98\% & 4\,s \\
    CatBoost            & \$220,406 & \$135,087 & 11.61\% & 10\,s \\
    Random forest       & \$228,130 & \$134,876 & 11.55\% & 14\,s \\
    Neural network      & \$247,526 & \$148,230 & 12.29\% & 49\,s \\
    Hedonic (OLS, log)  & \$306,564 & \$186,708 & 15.72\% & 4\,s \\
    \bottomrule
  \end{tabular}
\end{table}

The models also agree on what matters. Table~\ref{tab:shapmodels} shows
the SHAP share of the leading features in each model, computed on the
same 500 fold-5 test sales. Building floor area is first and distance to
the CBD second in all five models, and the coordinates are among the
top features in every model. The main difference is that the hedonic
model leans more on the SA2 income and employment statistics: a linear
model cannot learn sharp local patterns from latitude and longitude, so
it relies on smoother area-level summaries instead.

\begin{table}[t]
  \centering
  \caption{Share of total mean $|$SHAP$|$ (\%) for leading features,
  same 500 fold-5 test sales. RF = random forest, NN = neural
  network.}
  \label{tab:shapmodels}
  \begin{tabular}{@{}lrrrrr@{}}
    \toprule
    Feature & XGB & CatB & RF & NN & Hedonic \\
    \midrule
    Building floor area   & 21.0 & 21.8 & 31.6 & 16.8 & 21.0 \\
    Distance to CBD       & 15.5 & 14.6 & 18.9 & 16.6 & 11.0 \\
    Land area             & 12.7 & 14.0 & 13.6 & 11.5 & 3.8 \\
    Latitude + longitude  & 16.7 & 17.2 & 19.1 & 13.6 & 11.2 \\
    Deprivation index     & 6.9  & 8.1  & 3.7  & 7.0  & 3.4 \\
    Median income (SA2)   & 3.4  & 3.9  & 3.3  & 4.5  & 8.6 \\
    Building age          & 4.7  & 4.5  & 2.1  & 4.7  & 8.5 \\
    \bottomrule
  \end{tabular}
\end{table}

\subsection{Comparison with a hedonic model (RQ2)}\label{sec:res-hedonic}
The hedonic regression has the same inputs as XGBoost but is 4.7 MAPE
points less accurate (15.7\% vs 11.0\%), and its RMSE is 44\% higher
(\$306,564 vs \$212,678). It explains 72.9\% of the variance in log price
on fold~5's training data, so it is a reasonable model, but a straight
line in each feature cannot capture the non-linear and interacting
effects a tree model learns directly, such as a distance to the CBD that
matters a great deal near the city and very little far from it
(Fig.~\ref{fig:shap}).

In exchange, each hedonic coefficient has a direct, one-sentence
interpretation that holds for every house. On fold~5, for example, each
extra \$10,000 of median personal income in the SA2 is associated with
12.3\% higher prices, each kilometre further from the CBD with 1.2\%
lower prices, each step up the deprivation index with 1.3\% lower
prices, and a flood plain location with 1.8\% lower prices, all else
equal. SHAP gives a comparable dollar breakdown for the more accurate
model, but a feature's effect can change from one sale to another. This
is the accuracy/interpretability trade-off RQ2 asks about: the
machine-learning model keeps a per-sale explanation
(Section~\ref{sec:res-explain}) while cutting the hedonic model's MAPE by
about 30\%.

\subsection{Explaining a single prediction}\label{sec:res-explain}
Fig.~\ref{fig:house} shows the explanation for one held-out sale in fold
5's test period, a home sold in June 2026 for \$1,285,000. The XGBoost
model, trained only on sales up to January 2026, predicted \$1,266,131
(1.5\% below). Starting from the average prediction of \$1,151,033, the
neighbourhood raises the price by \$65,399, mostly through a low
deprivation index (2 of 10, +\$71,055); location adds \$41,119, with the
property's latitude (+\$85,033) outweighing its 16.9~km distance from
the CBD ($-$\$63,763); the house and land add \$26,184, as a large floor
area (172~m\textsuperscript{2}, +\$84,080) more than offsets a small
section (186~m\textsuperscript{2}, $-$\$94,139); and the 2026 market
timing subtracts \$19,070. Hazards change the price by only \$1,469. This
is the kind of reason an online valuation tool does not give: the buyer
can see that the price rests on the area and the size of the house, not
on the section.

\begin{figure}[t]
  \centering
  \figph{../figures/house_explanation.png}{4cm}{Waterfall explanation of
  one held-out sale.}
  \caption{SHAP explanation of one held-out sale (XGBoost, fold~5). Bars
  run from the average prediction to the model's prediction; red raises
  the price and blue lowers it. Lines mark the average, predicted and
  actual price.}
  \label{fig:house}
\end{figure}

\subsection{Comparison with online valuation estimates (RQ2)}
\ph{Online valuation benchmark results: see Section~\ref{sec:onlineval}
-- sample size, matched properties, and RMSE/MAE/MAPE on the matched
subset, alongside the augmented model's error on the same sales.}

%% ============= VI. LIMITATIONS, CONCLUSION, FUTURE WORK ==============
% About 1--1.25 pages. Rubric: Conclusion and Future Work (33 pts).
% Every claim must be supported in Section V.
\section{Limitations}\label{sec:limits}
Our price filter keeps only the middle 98\% of each fold's training
prices (Section~\ref{sec:split}), so our results describe the typical
Auckland home and may not transfer to the cheapest or most expensive
1\% of sales. Record linkage could not place a coordinate for every
sale: 61,244 arm's-length single-property sales were available, and
59,798 (97.6\%) were matched to a coordinate
(Section~\ref{sec:linkage}); if the unmatched 2.4\% are systematically
different (e.g.\ new subdivisions or non-standard addresses not yet in
LINZ), our results could be slightly optimistic or pessimistic for that
slice of the market, in a direction we have not separately quantified.

The neighbourhood features are 2023 Census and NZDep2023 values for the
whole SA2, so every home in an SA2 shares them, and they do not change
over the 2024--2026 sales period. They describe the area, not the
household that buys the home. Our attribution methods measure
association with the model's predictions, not causation: a high SHAP
value for the deprivation index says the model leans on that feature,
not that deprivation \emph{causes} a specific dollar change in price.
Because location and neighbourhood features overlap, the split of credit
between them (Section~\ref{sec:res-attribution}) depends on the method;
we report three methods for this reason, but no single number should be
read as the ``true'' value of one feature. Finally, we deliberately
fixed every model's hyperparameters for a fair comparison; a project
optimising purely for accuracy could tune each model and would likely
do somewhat better than every number reported here.

\section{Conclusion and Future Work}\label{sec:conclusion}
This paper asked whether Auckland-specific external data improves house
price prediction (RQ1), and how an explainable model built on it compares
with a traditional hedonic model and free online valuations (RQ2). For
RQ1, the answer is yes, and by a wide and statistically significant
margin: adding 16 external features cuts XGBoost's RMSE by 38.1\% and
MAPE by 8.4 percentage points, and a two-stage test confirms that the
external features explain 29.3\% of the error a property-only model
leaves. For the first half of RQ2, XGBoost is the most accurate of five
model families and far more accurate than a hedonic regression on the
same features (MAPE 11.0\% vs 15.7\%), while SHAP still explains each
prediction in dollars; the second half, comparison against free online
valuation estimates, is still being completed
(Section~\ref{sec:onlineval}).

The clearest practical finding is that most of the value of external
data in this project is the value of \emph{location}. Distance to the
CBD and the coordinates carry the largest share; area-level census and
deprivation statistics are useful but mostly repeat what location
already says; and flood, noise and development features add almost
nothing to a city-wide model. For a team building a similar model with a
limited data budget, accurate geocoding and a few well-chosen location
measures are worth more than many hand-engineered area features.

Future work should (1) finish and report the online valuation benchmark;
(2) tune each model and check whether the ranking of model families
holds; (3) test whether hazard features matter more in the parts of
Auckland where they are common, for example with location--feature
interaction methods such as GeoShapley~\cite{li2024geoshapley}, rather
than in the city-wide average reported here; and (4) extend the
record-linkage pipeline to recover the 2.4\% of sales that currently
have no matched coordinate, to check whether they change any of these
conclusions.

%% ======================== CONTRIBUTIONS =============================
% Not counted in the 12 pages. One or two specific sentences per member:
% coding, experiments, analysis, visualisation and writing.
\section*{Contributions}
\begin{description}
  \item[\ph{Member 1}:] \ph{Contribution.}
  \item[\ph{Member 2}:] \ph{Contribution.}
  \item[\ph{Member 3}:] \ph{Contribution.}
  \item[\ph{Member 4}:] \ph{Contribution.}
  \item[\ph{Member 5}:] \ph{Contribution.}
\end{description}



%% ========================== REFERENCES ==============================
\bibliographystyle{IEEEtran}
\bibliography{references.bib}

\end{document}
